\section{Introduction}

In the landscape of open-source software development, the number of GitHub stargazers has emerged
as a widely adopted proxy for project popularity and community engagement. Much like subscriber
counts on video platforms, stars reflect the perceived value of a repository and influence its
visibility, adoption, and long-term sustainability~\cite{borges2016}.

Predicting the star count of a repository from its activity metrics has practical value for
developers seeking to benchmark their projects, for investors evaluating open-source ecosystems,
and for researchers studying the social dynamics of software communities. However, the relationship
between repository attributes---such as fork counts, issue activity, and repository age---and star
accumulation is non-linear and multi-dimensional, making accurate prediction a non-trivial task.

In this project, we study which machine learning regression model achieves the highest prediction
accuracy for GitHub star counts. We collect a dataset of 1{,}000 open-source repositories with at
least 50 stars via the GitHub REST API, extract 14 activity-based features, and train five
regression models: Linear Regression, Ridge Regression, Random Forest, Gradient Boosting, and
Support Vector Regression (SVR). Models are evaluated using the coefficient of determination
($R^2$), and the best-performing model is deployed to a production server via an automated
Git~Hook pipeline.

The system is built on two Ubuntu virtual machines (VMs) provisioned on the SNIC Science Cloud
using OpenStack. Ansible orchestrates the configuration of both VMs, and Docker containers
encapsulate the application stack on the production server. This report describes the data
collection pipeline, the prediction system architecture, the comparative model results, and a
scalability analysis of the serving infrastructure.
